\section{Mixed propagation and unequal intervals}\label{c7:mixed}

Theorem~\ref{j:unequal} concerns two site lists whose sizes can differ.
Through \eqref{c7:eq-Z}, its proof reduces to estimating the signed
horizontal gas. Between the two suppression scales only the larger
list suppresses its species, so the interior operator from
Section~\ref{c7:compact} does not suffice there. We first bound the
operator for this intermediate region. We then extract the two exterior
power gaps and use the direction of the common-charge jumps in the
innermost region. The final subsection transfers the resulting
marked-bond estimate to planar polygons.

\subsection{The mixed spectral radius}

Fix the post-collision displacement $d=1/8$ and a sufficiently large
integer block size $b$, using Proposition~\ref{j:stability-eighth}.
On any fixed compact set of real angles, choose a common gauge with
quadratic slack.
For a constant real angle $\theta$, let $E_\theta$ be the homogeneous
operator allowing all particle types, and let $I_\theta$ allow only
$M,F_{\rm pair}$ in both species. Let $X_\theta$ allow all types in
species 1 and only $M,F_{\rm pair}$ in species 2. These are the
Hilbert--Schmidt operators constructed from the common gauge and
Bargmann tail in Section~\ref{c7:compact}. The empty-contour power
traces give
\[
 r(E_\theta)=r_E=e^{-\pi b/16}.
\]
The angles used below are $\theta=\pm\nu$, where $\nu=3\pi/4$.

\begin{proposition}[Mixed spectral bound]\label{c7:mixed-radius}
For every real $\theta$,
\begin{equation}\label{c7:eq-X}
 r(X_\theta)\le r_X^+,
 \qquad r_X^+=\exp\!\left[b\left(-\frac{5\pi}{96}
                              +\frac1{6\sqrt3}\right)\right].
\end{equation}
\end{proposition}

\begin{proof}
Use the regulator of Section~\ref{c7:markedgas} without $d_T$, with
constant angle $\theta$, and with homogeneous site counts $(0,m_*)$.
Here $m_*$ is an auxiliary site count. For a block count $N_*\ge2$,
keep $T=bN_*$ fixed until both site-count limits below have been taken.
The contour geometry gives a common period threshold $T_0$ independent
of the finite count $m_*$; the pre-collision domination constants may
depend on $m_*$. Enlarge $b$ so every such $T$ is at least $T_0$ for
both of the fixed displacement choices used below.
The unequal-list winding compensation is retained. Write
\[
 Q_{m_*,T}=\frac{m_*^2}{2}\int_T a_TD(1+g_dE)a_T.
\]
The two-species reorganization gives an exact finite-period identity
\begin{equation}\label{c7:two-species-comparison}
 \mathfrak C^{(0,m_*)}_T
 =c(T)e^{Q_{m_*,T}}\mathcal Z^{(2)}_{m_*,T},
\end{equation}
where $\mathcal Z^{(2)}_{m_*,T}$ is the gas at displacement $1/8$ with
these site powers. Its source is $(0,m_*a_T)$, so the quadratic source
contraction is exactly the species-2 diagonal entry $1+g_dE$.

The species-1 contour has no site poles. On a sufficiently small disk,
all of its positive-order Wick integrands are jointly analytic in their
contour variables, including at coincidences. Each such contour
integral vanishes, so this species contributes exactly 1 before the
Gaussian average. The remaining source and winding compensation depend
only on species 2. We can therefore evaluate the same regulator from
its original one-species marginal. Its nonzero Gaussian symbol is
$2\pi g_d/k$. Its constant mode is uniform modulo $2\pi$: the opposite
uniform mode remains uniform after addition of the independent common
Gaussian. Its winding $\ell\in8\mathbb Z$ has law
$e^{-4\pi\ell^2/T}/D_\ell$.

The one-species oscillator calculation uses the marginal before tilting.
After tilting, the height symbol is
\[
\begin{gathered}
 \mathsf h_* =\frac{g_dE^2}{1+g_dE}
 =E\bigl(1-(1+g_dE)^{-1}\bigr),\\
 1+g_dE=\frac{(1+t^3)^2(1+t^6)}{(1-t^4)(1-t^8)},\quad t=e^{-k},\\
 \frac{\mathsf h_*}{k}=\frac1{k^2}-\frac{41}{12}+O(k^2)\quad(k\downarrow0),
 \qquad \mathsf h_*=2e^{-3k}+O(e^{-4k})\quad(k\to\infty).
\end{gathered}
\]
Thus the one-species versions of $K^r,K_0$ use $\mathcal S=1$ and have
leading singularity $-2\log(1+s^2/9)$. The winding coefficients cancel:
\[
 \frac{41}{12}-4+\frac54-\frac23=0.
\]
These four terms come respectively from the covariance constant, the
winding law, the $H$ dress, and the modular scalar, in units
$\pi\ell^2/T$. The $\ell X^{(2)}$ term cancels the quadratic winding
dress as before. The inserted winding compensation cancels the remaining
$2\pi m_*\ell/T$ from the site-dependent $H$ constants.

The uniform constant mode unfolds the short heat sum. If $\mathsf X$ is the
charge-coordinate moment of the $H$ list, its integral is
\[
 \sqrt T\int_{\mathbb R}
 e^{-T\alpha^2/(4\pi)-\alpha\mathsf X}\frac{d\alpha}{2\pi}
 =e^{\pi\mathsf X^2/T}.
\]
It cancels the covariance factor $e^{-\pi\mathsf X^2/T}$, including the
Poisson factor $\sqrt T$. Hence the one-species gas has no state square
or initial-state sum. Its total particle charge must equal one
$\ell\in8\mathbb Z$. The leading linear interactions cancel in
modulus because $\mathcal S=1$, and the remaining real-position Berry
factors have modulus one. The local and residual rules are otherwise
those of Section~\ref{c7:boson}, with $\mathsf h_*$ in place of the
two-species height symbol. We use a fixed sufficiently small positive
post-collision displacement, reached from a sufficiently small negative
pre-collision displacement by the same contour deformation. The
resulting gas is denoted $\mathcal Z^{(*)}_{m_*,T}$. The second identity is
\begin{equation}\label{c7:one-species-comparison}
\begin{aligned}
 \mathfrak C^{(0,m_*)}_T
 &=c_*(T)e^{Q_{m_*,T}}\mathcal Z^{(*)}_{m_*,T},\\
 c_*(T)&=\frac{D_W}{D_\ell\operatorname{tr}\Gamma}
 \prod_{j\ge1}[1+g_d(2\pi j/T)E(2\pi j/T)]^{-1}.
\end{aligned}
\end{equation}
In particular, the deterministic exponential in the two identities
is the same.

\paragraph{Absolute summability in one species.}
We give the stability check needed to pass to the site-count limit in
\eqref{c7:one-species-comparison}. Use the normalized
$F=F_{\rm pair}/2$, and retain $M_0,X,f,g,l,j_0$ from
\eqref{j:stability-core}. The same parity decomposition and the relation
$(O_++O_-)/2=(P+M)/2$ apply. Use $\mathsf h_*$ in the height kernels;
the post-collision continuation coefficient is 2, and there is no
species off-diagonal subtraction. For the even core
$\mathcal A_*=M_0-\mathsf h_*XX^t$, direct reduction with
$u=e^{-k/2}$ gives
\[
 1-\mathsf h_*X^tM_0^{-1}X
 =\frac{(1-u^2)^2(1+u^4)^2}
 {3(1-u^2+u^4)(1-u^4+u^8)}>0.
\]
The rank-one criterion and $M_0>0$ give $\mathcal A_*>0$. Its endpoint
matrices are
\[
 \frac{\mathcal A_*}{k}\longrightarrow
 \begin{pmatrix}4&-2&-2\\-2&2&3/2\\-2&3/2&3/2\end{pmatrix}
 \quad(k\downarrow0),\qquad
 \mathcal A_*\longrightarrow\operatorname{diag}(1,1,1/4)
 \quad(k\to\infty).
\]
The first matrix has leading principal minors $4,4,1$. Continuity
therefore gives $\mathcal A_*\ge c\min(k,1)I$.

Here the displaced directions admit exact Schur complements. Put
$R_d=(U+L)/2$, $\delta_d=|d|$, $\Delta=\delta_d k$, and
\[
 \mathcal W_d=\frac{R_d-\cosh\Delta\,F}{\sinh\Delta}.
\]
Put $j_c=\cosh(k/2)/(2\cosh k)$ and $Y=-\sinh(3k/2)$.
The cross columns of $\mathcal W_d$ against the core, obtained from the
height lists, are
\[
\begin{aligned}
 \mathbf d_{\rm pre}
 &=\begin{pmatrix}-j_c\\-j_c\\-3/4-1/(4\cosh k)\end{pmatrix}
       +\mathsf h_*XY,\\
 \mathbf d_{\rm post}
 &=\begin{pmatrix}j_c\\j_c\\1/4+1/(4\cosh k)\end{pmatrix}
       -\mathsf h_*XY.
\end{aligned}
\]
In both cases its diagonal entry is
$\coth\Delta-g/2-\mathsf h_*Y^2$. Direct rational reduction gives
\[
\begin{aligned}
 g/2+\mathsf h_*Y^2+
 \mathbf d_{\rm pre}^t\mathcal A_*^{-1}\mathbf d_{\rm pre}
 &=\coth(k/2),\\
 g/2+\mathsf h_*Y^2+
 \mathbf d_{\rm post}^t\mathcal A_*^{-1}\mathbf d_{\rm post}&=E(k).
\end{aligned}
\]
Thus the Schur complements of $\mathcal W_d$ against the core are
\begin{equation}\label{c7:one-species-even-schur}
 S_{e,\mathcal W}^{\rm pre}=\coth\Delta-\coth(k/2),\qquad
 S_{e,\mathcal W}^{\rm post}=\coth\Delta-E(k).
\end{equation}
Here ``pre'' denotes the symmetrized absolute kernel at $d<0$, and
``post'' the continued kernel at $d>0$. In these entries the post continuation adds
$\sinh(2\Delta)$ to the $R_d,R_d$ entry and $\sinh\Delta$ to the
$R_d,F$ entry. Since
\[
 E(k)=\frac{1+t^3}{(1-t)(1+t^2)}<\frac{1+t}{1-t}=\coth(k/2),
\]
both Schurs are positive when $0<\delta_d<1/2$. We use this algebraic
range only at the sufficiently small contour displacements fixed above.

For the odd vectors $(O_+-O_-)/2$ and $(U-L)/2$, put
$\rho_k=\tanh(k/2)$ and $\tau_d=\tanh\Delta$. The Schur complements
against the first odd vector, divided by $\cosh^2\Delta$, are
\begin{equation}\label{c7:one-species-odd-schur}
\begin{aligned}
 \frac{S_o^{\rm pre}}{\cosh^2\Delta}
 &=\frac{(\rho_k-\tau_d)(\rho_k+3\tau_d)}{4\rho_k},\\
 \frac{S_o^{\rm post}}{\cosh^2\Delta}
 &=\frac{\rho_k L_o-J_o(2\tau_d+\tau_d^2/\rho_k)}{4D_o},
\end{aligned}
\end{equation}
where
\[
\begin{aligned}
 D_o&=2t^4-t^2+2,\\
 J_o&=2t^4-2t^3+t^2-2t+2,\\
 L_o&=6t^4+2t^3-5t^2+2t+6.
\end{aligned}
\]
The first odd pivot is $g>0$ pre-collision and
$g+\mathsf h_*\sinh^2k>0$ post-collision. Also $\tau_d<\rho_k$,
\[
 D_o=2(t^2-1/4)^2+15/8>0,\quad
 J_o=2(1-t)^2(t^2+t+1)+t^2>0,\quad
 L_o-3J_o=8t(t^2-t+1)>0.
\]
The post expression in \eqref{c7:one-species-odd-schur} decreases with
$\tau_d$; at $\tau_d=\rho_k$ it is
$2\rho_k t(t^2-t+1)/D_o>0$. This proves both odd signs.

Restoring the original $R_d$ coordinate multiplies the even Schur in
\eqref{c7:one-species-even-schur} by $\sinh^2\Delta$. For fixed
$0<\delta_d<1/2$, the limits of the original Schurs divided by $k$ at
zero are
\[
\begin{array}{c|cc}
 &\mathrm{pre}&\mathrm{post}\\\hline
 S_e/k&\delta_d(1-2\delta_d)&\delta_d(1-\delta_d)\\
 S_o/k&(1-2\delta_d)(1+6\delta_d)/8&(11-4\delta_d-4\delta_d^2)/24.
\end{array}
\]
The original cross entries are $O(k)$ there. At infinity the added
Schurs tend to $1/2$ and the cross entries tend to zero. For the extreme
post height terms the required cancellations are explicitly
\[
 -\tfrac12e^{2\Delta}+\sinh(2\Delta)=-\tfrac12e^{-2\Delta},\qquad
 -\tfrac12e^{\Delta}+\sinh\Delta=-\tfrac12e^{-\Delta};
\]
the remaining high-frequency terms decay because $\delta_d<1/2$.
Together with the core bounds, these limits give a lower bound
$c_d\min(k,1)$ on the count-respecting quotient for each fixed nonzero
$d$ used here. Its dependence on $d$ matters: the count functional has
value $1/2$ on normalized $F$ and value 1 on $R_d$, so the extra null
direction $R_0-F$ at $d=0$ is visible to count.

The physical high-frequency matrix is entrywise nonnegative, with the
same endpoint coincidences as before and normalized $F,F$ entry $1/4$.
Split off this matrix times $\tanh(e k)$ with a sufficiently small
fixed $e>0$. The remaining positive symbol divided by $k$ is regular
at zero; there is no common-charge subtraction in one species. The
periodized argument of Section~\ref{c7:stability} therefore controls
the sum of squared box counts.

For precision, its pre-collision base product is the one-species
counterpart of \eqref{j:pre-weight}. Replace every $H$ atom $(a,r)$ by
$(a,r/2),(-a,r/2)$ in both the local rule
\eqref{c7:local-factor-rule} and the forward rule
\eqref{c7:forward-residual}, using the scalar $K_0,K^r$ determined by
$\mathsf h_*$. This symmetrizes the internal $H$ self-list as well.
These local factors are finite on the chosen negative contours, since
their extreme symmetrized height difference is $3+2d<3$.
Keep the short and outer lists at the same negative displacement, take
the moduli, and multiply all local factors and all forward factors,
including nonzero self-images. At the real angle used here, the heat
modulus replaces $\mathsf X$ by $\Re\mathsf X$. The unfolded integral above cancels
$e^{-\pi(\Re\mathsf X)^2/T}$, including its Poisson normalization. The scalar state
coefficient $1-\mathcal S$ is zero, so this base product has state factor
one and no initial-state sum. The matrix bound just proved applies to
this product. The remaining deterministic pre-collision factors have
a finite bound per particle at each fixed $(m_*,T,d)$, and their linear
count cost is absorbed by the quadratic bound. Each configuration fixes
at most one permitted winding. This proves absolute summability for the
actual pre-collision Gaussian moment and the post gas, justifying
\eqref{c7:one-species-comparison}. On the post contours it also provides
domination independent of $m_*$ at fixed $T$.

\paragraph{The fixed-period site-count limit.}
In either gas, the site multipliers have modulus at most one on its
chosen post contours. At fixed $T$, as $m_*\to\infty$, all suppressed
species-2 types disappear and $M,F_{\rm pair}$ retain multiplier one.
The preceding count bound and the two-species stability bound justify
dominated convergence. Thus
\[
 \lim_{m_*\to\infty}\mathcal Z^{(2)}_{m_*,bN_*}
 =\det(I-A_0^{N_*})\operatorname{tr}X_\theta^{N_*}.
\]
We next bound the corresponding one-species limit.

For its remaining types $M,F_{\rm pair}$, the negative logarithms of
the residual pair magnitudes have Fourier symbols, in units $2\pi/k$,
\[
\begin{array}{c|c}
 MM&(1-t^4)(1-t+t^2)/(1+t^6)\\
 FF&(1-t^4)(1+t+t^2)/(1+t^6)\\
 MF&t^{1/2}(1-t^3)(1+t^2)/(1+t^6).
\end{array}
\]
These are the $M,F$ entries of $M_0-\mathsf h_*XX^t$ after restoring
$F_{\rm pair}=2F$. To invert them, put
$D_x=\cosh(\pi x/6)$ and $R_x(a)=(D_x-a)/(D_x+a)$. The hyperbolic log
transform gives
\[
\begin{aligned}
 P_{MM}(x)&=R_x(1)R_x(1/2)/R_x(\sqrt3/2),\\
 P_{FF}(x)&=R_x(1)R_x(1/2)R_x(\sqrt3/2),\\
 P_{MF}(x)&=R_x(\cos(\pi/12))R_x(\cos(5\pi/12)).
\end{aligned}
\]
Indeed $-\log R_x(\cos(\pi h/6))$ has symbol
$(t^h-t^{6-h})/(1+t^6)$ in these units for $0\le h\le3$. All three pair magnitudes are
at most one for $x>0$. For $P_{MM}$ use
$R_x(1)\le R_x(\sqrt3/2)$ and $R_x(1/2)\le1$; the other two bounds are
immediate. This includes every forward image factor.

The absolute local constants for $i=M,F_{\rm pair}$ satisfy
\[
 z_i=\lim_{x\downarrow0}\frac{\sqrt{P_{ii}(x)}}{2\pi x}.
\]
For $M$, the field prefactor has magnitude $1/\sqrt8$, and its internal
height exponent is half the coincident two-body contraction. The short
and outer zeros give $(2\pi x)(\pi x/4)$, proving this rule. The linear
part of $K^r$ has no contribution to the limiting real part. For
$F_{\rm pair}$, set
$\rho_{\rm col}=\lim_{\varepsilon\to0}|\varepsilon|^2
 e^{\Re K_0(3i+\varepsilon)}$. Its collision residue has magnitude
$(4\pi^2/8)e^{K_0(0)}\rho_{\rm col}$: this uses the short-pair zero and
the unit outer magnitude at height difference one. The leading
two-body factor is
\[
 (2\pi x)^4(\pi x/4)^2e^{2K_0(0)}\rho_{\rm col}^2/x^4,
\]
which gives the same local rule. Consequently
\[
 z_M=\frac{2+\sqrt3}{24\sqrt3},\qquad
 z_F=\frac{2-\sqrt3}{24\sqrt3},\qquad
 z_M+z_F=\frac1{6\sqrt3}.
\]
Dropping the pair factors and the winding constraint only increases
the absolute sum. The indistinguishable-particle factorials then give
\[
 \left|\lim_{m_*\to\infty}\mathcal Z^{(*)}_{m_*,T}\right|
 \le e^{T(z_M+z_F)}=e^{T/(6\sqrt3)}.
\]
Divide \eqref{c7:two-species-comparison} and
\eqref{c7:one-species-comparison} by their common nonzero exponential,
and take these two limits at the fixed period. We obtain
\[
 |\operatorname{tr}X_\theta^{N_*}|
 \le |\det(I-A_0^{N_*})|^{-1}
      \left|\frac{c_*(bN_*)}{c(bN_*)}\right|
      e^{bN_*/(6\sqrt3)}.
\]
Using $\int_0^\infty\log(1+e^{-ak})\,dk=\pi^2/(12a)$ and the
corresponding minus-sign integral from Section~\ref{c7:boson}, the
determinant product in $c_*$ has rate $-19\pi/288$; together with
$11\pi/144$ from $D_W/\operatorname{tr}\Gamma$, this gives
$\log c_*(T)/T\to\pi/96$. Since $\log c(T)/T\to\pi/16$,
\[
 \frac1T\log\frac{c_*(T)}{c(T)}\longrightarrow-\frac{5\pi}{96}.
\]
Now let $N_*\to\infty$. The determinant tends to one, and the
power-trace spectral-radius criterion from Section~\ref{c7:compact}
gives \eqref{c7:eq-X}. The fixed-site physical limit has not entered
this comparison.
\end{proof}

\subsection{A coefficient lemma for two exterior gaps}

The fixed-list limit in \eqref{c7:eq-Z} controls a trace with two long
exterior portions, one at each limiting angle. The following version
of Lemma~\ref{j:trace-extraction} records the frequency collision that
occurs when both portions have length $2N-L$.

\begin{lemma}[Two exterior gaps]\label{c7:two-gap}
Let $E_+,E_-$ be Hilbert--Schmidt operators on a complex Hilbert space,
each with spectral radius one. Write
$\mathcal P_\pm=\{\zeta\in\sigma(E_\pm):|\zeta|=1\}$ for their finite
nonempty peripheral spectra. For $\zeta\in\mathcal P_\pm$, let
$P_\zeta^\pm$ be the Riesz projection,
$J_\zeta^\pm=(E_\pm-\zeta I)P_\zeta^\pm$, and choose
$q_\zeta^\pm\ge1$ to be its nilpotent order on
$\operatorname{ran}P_\zeta^\pm$.
Let $C$ be trace class. For $i=1,2,3,4$, let $\mathscr W_i$ be a
countable family of bounded operators with integer lengths $\ell(W)\ge0$
such that, for fixed $\varepsilon,B>0$,
\[
 \sum_{W\in\mathscr W_i}e^{\varepsilon\ell(W)}\|W\|\le B.
\]
Fix an integer $L\ge0$. For integers $N\ge1$ with $g=2N-L\ge0$, put
\[
\begin{gathered}
 F_{N,L}=\sum_{\substack{W_i\in\mathscr W_i\\
                \ell_+\le g,\ \ell_-\le g}}
 \operatorname{tr}\bigl(CW_1E_+^{g-\ell_+}W_2W_3
                           E_-^{g-\ell_-}W_4\bigr),\\
 \ell_+=\ell(W_1)+\ell(W_2),\quad
       \ell_-=\ell(W_3)+\ell(W_4).
\end{gathered}
\]
Suppose $F_{N,L}-R\operatorname{tr}E_+^{4N}\to0$ for a scalar $R$ as
$N\to\infty$ at this fixed $L$.
For either sign define
\[
 \mathcal B_\zeta^\pm(z,\ell)=
 \sum_{a=0}^{q_\zeta^\pm-1}
 \binom{z-\ell}{a}\zeta^{-a}(J_\zeta^\pm)^aP_\zeta^\pm.
\]
The following series is coefficientwise absolutely convergent and defines
a polynomial:
\[
\begin{aligned}
 \mathscr P_{\zeta,\eta}(z)={}&
 \sum_{W_1,W_2,W_3,W_4}\zeta^{-\ell_+}\eta^{-\ell_-}\\
 &\quad\cdot\operatorname{tr}\bigl(CW_1\mathcal B_\zeta^+(z,\ell_+)W_2W_3
                              \mathcal B_\eta^-(z,\ell_-)W_4\bigr).
\end{aligned}
\]
Then for every frequency $\beta$ in
$\{\alpha^4:\alpha\in\mathcal P_+\}\cup
 \{(\zeta\eta)^2:\zeta\in\mathcal P_+,\eta\in\mathcal P_-\}$,
\begin{equation}\label{c7:two-gap-frequency}
\begin{aligned}
 R\sum_{\substack{\alpha\in\mathcal P_+\\\alpha^4=\beta}}
       \operatorname{rank}P_\alpha^+
 &=\sum_{\substack{\zeta\in\mathcal P_+,\eta\in\mathcal P_-\\
                   (\zeta\eta)^2=\beta}}
       (\zeta\eta)^{-L}\mathscr P_{\zeta,\eta}(-L).
\end{aligned}
\end{equation}
In particular
\[
 |R|\le K(1+L)^h\|C\|_1,
 \qquad h=\max_{\zeta,\eta}(q_\zeta^++q_\eta^--2),
\]
where $K$ depends only on $E_\pm$, $\varepsilon$, and $B$, and is
independent of $C,L,R$.
\end{lemma}

\begin{proof}
The peripheral expansions from Lemma~\ref{j:trace-extraction} give,
for either operator, a finite sum of its Riesz nilpotent polynomials
plus an exponentially small remainder in operator norm. Its powers
have a fixed polynomial bound. The exponential word moments make the
displayed series for $\mathscr P_{\zeta,\eta}$ coefficientwise
absolutely convergent: a coefficient of $\binom{z-\ell}{a}$ is at most
$C_a(1+\ell)^a$. Thus every coefficient has magnitude at most
$K\|C\|_1$, and the degree is at most $h$.

For word quadruples with both $\ell_+,\ell_-\le g/2$, replace both
powers in $F_{N,L}$ by their peripheral expansions. Their total error
is exponentially small in $g$ times a polynomial, by the word bounds.
If either total word length exceeds $g/2$, its exponential tail absorbs
the polynomial power bounds. The same estimate permits extending the
peripheral sums to all quadruples. Hence, at fixed $L$,
\[
 F_{N,L}=\sum_{\zeta\in\mathcal P_+,\eta\in\mathcal P_-}
           (\zeta\eta)^g\mathscr P_{\zeta,\eta}(g)+o(1).
\]
The bare trace is
$\sum_{\alpha\in\mathcal P_+}d_\alpha\alpha^{4N}+o(1)$, with
$d_\alpha=\operatorname{rank}P_\alpha^+$; positive nilpotent powers
have trace zero. Rewrite $g=2N-L$ and group the frequencies in $N$.
Uniqueness of unit-circle exponential polynomials, proved by Ces\`aro
extraction in Lemma~\ref{j:trace-extraction}, gives the polynomial
identity
\[
 \sum_{(\zeta\eta)^2=\beta}(\zeta\eta)^{-L}
       \mathscr P_{\zeta,\eta}(2z-L)
 \equiv R\sum_{\alpha^4=\beta}d_\alpha.
\]
Its value at $z=0$ is \eqref{c7:two-gap-frequency}. Choose a frequency
$\beta=\alpha_0^4$ from $\mathcal P_+$. The sum on the left side of
\eqref{c7:two-gap-frequency} multiplying $R$ is a positive integer.
Since $|(\zeta\eta)^{-L}|=1$ and the polynomial coefficients are
bounded as above, evaluation at $-L$ proves the asserted bound.
\end{proof}

\subsection{Proof of the unequal-interval estimate}

\begin{proof}[Proof of Theorem~\ref{j:unequal}]
Fix the bound on the count discrepancies. We first use distinct generic
imaginary site heights in the single neighborhood supplied by
Proposition~\ref{j:shifted-suppression}, with $w_e=-if$ and with the
site heights together with the endpoint height $-f$ having total spread
strictly below one. All constants below are uniform over these nearby heights,
$0\le f\le0.6$, and the list orders. The interval sizes are $m+1,n+1$,
where $n\ge m$ and $M=m+n+2$. Formula \eqref{c7:eq-Z} applies to the
continued matrix quantity $H^{\rm flip}$ on each such fixed pair of
lists.

Use one homogeneous gauge with quadratic slack for all kernels in this
proof. The bounded angle profiles and bounded extra dress logarithms
cost only linear terms in the local counts and common charge, which the
quadratic gauge absorbs uniformly. The Bargmann tail is independent of
these profiles. Normalize every block by $r_E$. The resulting kernels
have uniform Hilbert--Schmidt bounds. A perturbation of size $\delta$
in the angle and per-particle logarithms changes a kernel by
$O(\delta)$ in Hilbert--Schmidt norm; polynomial count and charge
factors from exponentiating the difference are absorbed by the same
quadratic bound. For suppressed multipliers of modulus at most one,
the corresponding estimate uses their differences directly, so it also
applies at their zeros.

\paragraph{Removing the distant periodic copies.}
Take $T=4Nb$ and let $N\to\infty$ while the two lists are fixed.
On blocks $j=-N,\ldots,N-1$, use the single origin kink
$D_e(x-w_e)$, its single localized dress corrections from
\eqref{c7:eq-M}, and the unperiodized products of all site factors.
On blocks $j=N,\ldots,3N-1$, use the opposite kink translated to
$2Nb$, negate those angle and correction logs, and omit all site
products. Exponential image tails show that each replacement of a
periodized kernel costs $O_{m,n}(e^{-cN})$ at these fixed lists.

Put
\[
 L=\left\lceil\frac{\log n}{\kappa_1b}\right\rceil,
 \qquad \kappa_1=\pi/3,
\]
and let $\mathcal K_c$ be the ordered product of the normalized origin
profile kernels on $-L\le j<L$. For large $N$ it is followed by a
positive-angle portion of length $g=2N-L$ and a negative-angle portion
of the same length. Their homogeneous operators are
$E_+'=E_{\nu}/r_E$ and $E_-'=E_{-\nu}/r_E$.

On each portion, the deviations from the homogeneous operator decay
as $Ce^{-ck}$ in block distance $k$ from its two ends. At the end next
to $\mathcal K_c$, the site log tail is bounded because
$n e^{-\kappa_1bL}\le1$; the kink corrections also decay there. At the
other end only the opposite kink contributes. These constants are
uniform in the lists. Expand each end in subsets of its deviations,
ending a word at its last chosen deviation. Pair each intervening
polynomially bounded power of $E_\pm'$ with the farther chosen index.
The exponential deviation bounds then give four word families with
\[
 \sum_W e^{c'\ell(W)}\|W\|\le C
\]
for fixed $c'>0$, uniformly in the lists. The same expansion bounds
every subinterval product polynomially in $N$ at fixed lists.
Consequently the $O_{m,n}(e^{-cN})$ periodization errors have total
trace contribution $o(1)$; keep two Hilbert--Schmidt factors and use
the polynomial bounds on the intervening products.

Let $F_N^{\rm per}$ be the normalized trace with the periodized
kernels. The exact trace formulas, including their common Bargmann
determinant, say
\[
 F_N^{\rm per}=c(4Nb)Z^L_{4Nb}\operatorname{tr}(E_+')^{4N}.
\]
For these fixed lists, \eqref{c7:eq-Z} supplies
$R=\lim_{N\to\infty}c(4Nb)Z^L_{4Nb}$. The bare power traces are bounded
by their peripheral expansion even when $E_+'$ has Jordan blocks.
It follows that $F_N^{\rm per}-R\operatorname{tr}(E_+')^{4N}\to0$.
After the periodization replacement the same limit holds for the
profile trace. Its four anchored words occupy portions of lengths
$N-L$ and $N$ in each gap. Extending those separate cutoffs to the
word sums of Lemma~\ref{c7:two-gap} adds $o(1)$ at fixed $L$: every
extra word has length at least $N-L$, and its exponential tail absorbs
the polynomial powers. The lemma applies with $C=\mathcal K_c$.
We have proved the uniform estimate
\begin{equation}\label{c7:eq-G}
 \left|\lim_{N\to\infty}c(4Nb)Z^L_{4Nb}\right|
 \le C(1+L)^h\|\mathcal K_c\|_1.
\end{equation}
The dependence of the discarded periodization errors on the fixed
lists has no effect on the uniform constant in this coefficient bound.

\paragraph{The central charge estimate.}
Put
\[
 J=\left\lfloor\frac{\log m}{\kappa_1b}\right\rfloor,
 \qquad W=\lfloor\gamma L\rfloor
\]
for a fixed small buffer fraction $\gamma>0$. Radii here count full
blocks from the origin. On each side, ignore
buffers around radii $0,J,L$. The remaining inner comparison has
length $(J-2W)_+$ and the remaining mixed comparison has length
$(L-J-2W)_+$, up to bounded endpoint changes. Thus the ignored length
relative to $J$ and $L-J$ is $O(W+1)$ even when intervals overlap or
are empty. Buffer products cost at most $C^{O(W+1)}$; reserve one
Hilbert--Schmidt factor at each end of the whole product.

At radii between $J+W$ and $L-W$, the small-list multipliers differ
from one by $O(e^{-cW})$, and the large-list suppressed multipliers are
superexponentially small in $W$. The angle and localized dress differ
from their limiting values by $O(e^{-cW})$. Hence these normalized
kernels are $O(e^{-cW})$ from $X_{-\nu}/r_E$ on the left and
$X_{\nu}/r_E$ on the right.

At radii below $J-W$, omit both species' suppressed labels, including
through the center buffer. The total error in the central product is
at most
\[
 L C^{2L}\exp(-c e^{cW}),
\]
with zero error if this inner interval is empty. This follows directly
from suppression at that distance inside the smaller scale and the
crude uniform bound for every other kernel. Between radii $W$ and
$J-W$, the remaining normalized kernels are $O(e^{-cW})$ from
$I_{-\nu}/r_E$ and $I_{\nu}/r_E$ on the two sides.

Only after this omission do we restrict charge sectors. A spatial
successor satisfies $p'=p-q$, and the remaining particles have
$q_M=-1$ and $q_{F_{\rm pair}}=-2$ in their species. Therefore
$k=p_1+p_2$ is nondecreasing along the whole restricted inner interval.
This is the spatial direction; the operator acts from successor to
predecessor. Partition the product by the sign of $k$ at its midpoint.
For the left and right angles define
$d_-(k)=|k-3/4|$ and $d_+(k)=|k+3/4|$. Monotonicity permits the
following comparison ranges $\mathcal R_-,\mathcal R_+$; the columns
$k_-^*,k_+^*$ give the spectral maximizers in those ranges:
\[
\begin{array}{c|cc|cc|c}
 \text{midpoint }k&\mathcal R_-&k_-^*&\mathcal R_+&k_+^*&
 d_-(k_-^*)^2+d_+(k_+^*)^2\\\hline
 0&\{k\le0\}&0&\{k\ge0\}&0&9/8\\
 >0&\mathbb Z&1&\{k\ge1\}&1&25/8\\
 <0&\{k\le-1\}&-1&\mathbb Z&-1&25/8
\end{array}
\]
For any range $\mathcal R$ in the table, the
compression of $I_\theta$ to its common-charge fibers has radius
\[
 \max_{k\in\mathcal R}u_k(\theta),\qquad
 u_k(\theta)=\exp\!\left[-\frac{4\pi b}{9}
                              (k+\theta/\pi)^2\right].
\]
Every closed trace in a compression contains only empty blocks, because
all allowed particles have negative common charge. The empty-block power traces
therefore give this radius, as in Section~\ref{c7:compact}. This argument
works at the real angles $\pm\nu$ and needs no semisimplicity assertion.
The last column of the table shows that every product of the applicable
left and right radii is at most
\[
 u_c^2=e^{-(4\pi b/9)(9/8)},\qquad u_c=e^{-\pi b/4}.
\]

For each of the finitely many comparison operators, a spectral-radius
upper bound $r>0$ gives $\|B^s\|\le C_\delta(re^\delta)^s$ for every
fixed $\delta>0$. Expanding the $O(e^{-cW})$ kernel deviations by
subsets gives the same bound with an arbitrarily small further spectral
slack once $W$ is large. All constants are uniform because the
comparison operators are the fixed $X_{\pm\nu}/r_E$ and the finitely
many compressed $I_{\pm\nu}/r_E$. Combine these bounds with the table,
Proposition~\ref{c7:mixed-radius}, and the buffer costs. For fixed
$\gamma,\delta>0$ the result is
\[
\begin{split}
 \|\mathcal K_c\|_1\le{}&C_{\gamma,\delta}
 e^{C_b\gamma L+2\delta L}
 (u_c/r_E)^{2J}(r_X^+/r_E)^{2(L-J)}\\
 &+LC^{2L}e^{-c e^{cW}}.
\end{split}
\]
The additive error is negligible even relative to an exponential in
$-L$ for each fixed $\gamma>0$. Since
$J=\log m/(\kappa_1b)+O(1)$ and
$L-J=\log(n/m)/(\kappa_1b)+O(1)$, choosing arbitrarily small fixed
$\gamma,\delta$ proves that for every $\eta>0$
\[
 \|\mathcal K_c\|_1\le C_\eta n^\eta
 (u_c/r_E)^{2\log m/(\kappa_1b)}
 (r_X^+/r_E)^{2\log(n/m)/(\kappa_1b)}
\]
whenever $n\ge m$ are sufficiently large. The threshold can be imposed
on $m$, and may depend on the fixed count-discrepancy bound. This also
covers short or absent comparison segments through the stated buffer
allowance.

The polynomial factor in \eqref{c7:eq-G} is absorbed by an arbitrarily
small power of $n$. Combining that equation with the uniform version
of \eqref{c7:eq-Z}, using $M\asymp n$, gives the desired power because
\[
\begin{aligned}
 -\frac5{24}+\frac{2\log(u_c/r_E)}{\kappa_1b}&=-\frac43,\\
 -\frac5{24}+\frac{2\log(r_X^+/r_E)}{\kappa_1b}
 &=\frac1{\pi\sqrt3}-\frac7{48}<\frac13.
\end{aligned}
\]
For each prescribed $\epsilon>0$, choose the induced Pfaffian, buffer,
spectral, and polynomial exponents with total at most $\epsilon$. This proves a
bound with constant $C_\epsilon$ uniform over the nearby generic
heights and orders. Now approach the physical spread-one heights for
each fixed pair of lists. Proposition~\ref{j:marked-physical} gives
regularity of $H^{\rm flip}$ and its identification with
$H_{\rm lg}=H^{\rm poly}$ there. No regulator limit at the boundary is
needed. We conclude, in the uniform notation of the introduction,
\begin{equation}\label{c7:eq-U}
 H_{\rm lg}\le n^{o(1)}m^{-4/3}(n/m)^{1/3}.
\end{equation}
This is the assertion of Theorem~\ref{j:unequal}, with constants and
threshold allowed to depend on the fixed discrepancy bound.
\end{proof}

\subsection{From selected bond pairs to planar polygons}

We finish the length-square estimate for planar polygons modulo lattice
translations:
\begin{equation}\label{c7:eq-Pol}
 \sum_{\lambda:\,\operatorname{diam}\lambda\le H}
 x^{|\lambda|}|\lambda|^2\le H^{2/3+o(1)}.
\end{equation}
The argument first removes polygons with too few locally available bond
directions. For the remaining polygons it selects many bond pairs and
constructs, for each pair, a cylinder satisfying every hypothesis of
Theorem~\ref{j:unequal}.

Fix a large integer $K$. A vertex of a polygon is called good if the
cyclic window of $2K$ path edges centered at that vertex uses bonds
crossing tessellation sides
of all three orientations. Write $N=|\lambda|$ for the polygon length.

\begin{lemma}[Summability of polygons with few good vertices]
\label{c7:bad-windows}
For a sufficiently large fixed $K$,
\[
 \sum_{\lambda:\,\#\{\text{good vertices of }\lambda\}<|\lambda|/2}
 x^{|\lambda|}|\lambda|^2<\infty,
\]
where polygons are counted modulo translations.
\end{lemma}

\begin{proof}
For large $N$, at least $N/2$ centers have a window omitting one
orientation. Greedily select disjoint cyclic radius-$K$ windows at
such centers. A selected window excludes at most $4K+1$ possible
centers, and contains $2K$ path edges. The number of covered edges,
after discarding the possible window crossing a chosen root, is at least
\[
 2K\left(\frac{N/2}{4K+1}-1\right)-2K\ge \frac N7
\]
for $N$ sufficiently large at fixed $K$.

Count the resulting rooted oriented paths. There are at most $N/(2K)$
selected window starts. Specifying those starts, one missing orientation
and the first step of each window costs at most
\[
 \sum_{j\le N/(2K)}\binom Nj6^j
 \le \exp(\varepsilon_K N+o(N)),\qquad \varepsilon_K\longrightarrow0
 \quad(K\to\infty).
\]
Inside a selected window continuation is forced: after one of the
three bond orientations is forbidden, the allowed graph has degree
two, and a nonbacktracking path has one next edge. Outside the windows
there are at most two choices per edge. Up to polynomial factors from
rooting and closing, the number of paths is therefore at most
$e^{\varepsilon_KN+o(N)}2^{6N/7}$. Since
\[
 x2^{6/7}=(2x)2^{-1/7}<1,
\]
choose $K$ so the entropy factor preserves a strict exponential decay.
The extra factor $N^2$ and the rooting factors are then summable. The
finitely many shorter lengths contribute a finite amount.
\end{proof}

\begin{lemma}[Selected pairs and permitted cylinders]\label{c7:selected-pair}
Let $\lambda$ be a planar polygon of diameter at most $H$, and let
$v_1,v_2$ be good vertices. Put $r_K=(K-1/2)/\sqrt3$. If
$|v_2-v_1|>4\sqrt{19}\,r_K$, one can assign used bonds $e_1,e_2$ within
their respective windows so that the following holds after a lattice
rotation or reflection. Both bonds cross horizontal tessellation sides,
and their midpoint displacement is
\[
 \delta=a+b\tau,\qquad a\ge1,\quad b\ge0,\quad
 f=\frac{b}{a+b}<0.6.
\]
There is a permitted cut with these two horizontal marked sites and,
after a common rescaling, physical logarithmic heights $-f,1-f$. Its
interval sizes are $m+1,n+1$, its period has an even number
$M=m+n+2$ of sites, and its unmarked count discrepancies from the
proportions $1-f,f$ are at most $f$. It has $n\ge m$,
$m\asymp|\delta|$, and $n\asymp H$
above a fixed distance cutoff. Projection to this cylinder is injective
on the polygons of diameter at most $H$ containing the fixed ordered
bond pair, when their first bond is fixed in the plane.
\end{lemma}

\begin{proof}
Choose the tessellation ray closest to $v_2-v_1$. Its angle from that
displacement is at most $\pi/6$. Goodness supplies, in each window, a
used bond crossing a side parallel to this ray. Its midpoint is within
$r_K$ of the corresponding vertex. Let $\delta$ be the displacement
of these midpoints. Set $\theta_* =\arg(2+3\tau)$. Elementary geometry
gives
\[
 \sin(\theta_*-\pi/6)=\frac1{2\sqrt{19}}.
\]
Since $|\delta-(v_2-v_1)|\le2r_K$, the change of angle is at most
$\arcsin(2r_K/|v_2-v_1|)$. The stated distance cutoff therefore makes
the angle of $\delta$ from the selected ray strictly less than
$\theta_*$. Rotate the selected ray to the positive horizontal direction,
then reflect if necessary to make that angle nonnegative.
Midpoints of horizontal-type bonds form one coset of
$\mathbb Z+\mathbb Z\tau$. Thus $\delta=a+b\tau$ with integers
$a\ge1,b\ge0$. Moreover
\[
 \tan\arg\delta=\frac{\sqrt3 f}{2-f},\qquad f=\frac b{a+b},
\]
and the value at $f=3/5$ is $\tan\theta_*$, proving $f<0.6$.

Put $s=a+b$ and $d=|\delta|$. Then
\[
 d=s\sqrt{1-f+f^2},\qquad \frac{\sqrt3}{2}s\le d\le s.
\]
Both midpoints lie in the convex hull of the polygon vertices, whose
diameter is also at most $H$, so $d\le H$. Construct a word of
$a$ horizontal and $b$ shifted $\tau$ steps, beginning with the
horizontal side crossed by $e_1$. Its next copy begins with the side
crossed by $e_2$. Repeat this word $2L$ times, where
\[
 L=\left\lfloor\frac{3H}{2d}\right\rfloor+1.
\]
For a lattice point $p+q\tau$, the transverse coordinate $p+q$
increases by one along either allowed cut step and is unchanged by the
time step $\tau^2=\tau-1$. The translated rows therefore tile the
plane and the cut is permitted.

The period vector and site count are $P=2L\delta$ and $M=2Ls$.
Thus $|P|>3H$, $M$ is even, and
\[
 3H<M\le\frac{10}{\sqrt3}H.
\]
The upper bound uses $L\le3H/(2d)+1$, $s/d\le2/\sqrt3$, and $d\le H$.
The two interval sizes and their exact unmarked counts are
\[
\begin{gathered}
 m=s-1,\qquad n=(2L-1)s-1,\\
 (a-1,b)=\bigl(m(1-f)-f,\ mf+f\bigr),\\
 ((2L-1)a-1,(2L-1)b)
   =\bigl(n(1-f)-f,\ nf+f\bigr).
\end{gathered}
\]
The two marked sites are horizontal, and every count discrepancy is
at most $f<0.6$. Apply a common rescaling of the physical site variables
so the horizontal logarithmic height is $-f$. The shifted-to-horizontal
ratio $q$ then makes the shifted height $1-f$, as required by
Theorem~\ref{j:unequal}. Since $d\le H$, one has $L\ge2$, hence $n\ge m$.
The displayed comparisons give $m\asymp d$ above a fixed cutoff and
$n\asymp H$.

Finally, no two points of a diameter-$H$ polygon differ by a nonzero
multiple of $P$, because $|P|>3H$. Its projection is therefore an
embedded cylinder polygon through the two marked bonds. The lift of a
connected projected polygon containing the fixed first plane bond is
unique. Two plane polygons containing that fixed ordered pair cannot
have the same projected polygon. This proves the claimed injection.
\end{proof}

Choose $K$ as in Lemma~\ref{c7:bad-windows}. For each prescribed
exponent slack $\epsilon>0$, choose a fixed distance cutoff large enough
both for Lemma~\ref{c7:selected-pair} and for its resulting
$m\ge m_\epsilon$ in Theorem~\ref{j:unequal}. In a polygon with at least
$N/2$ good vertices, lattice packing excludes only $O(N)$ ordered pairs
within this fixed cutoff. Thus for sufficiently large $N$ the retained
pairs number at least $cN^2$. Choosing a used bond in each window loses
only a bounded multiplicity: each selected bond can be assigned by only
$O(K)$ vertices along the polygon. The finitely many shorter lengths
contribute a bounded total modulo translations.

Fix the first bond at one translation representative of each bond type.
In a dyadic displacement annulus of radius $d$, there are $O(d^2)$
possible second bonds. For each selected pair among these possibilities,
Lemma~\ref{c7:selected-pair}
embeds the diameter-truncated plane polygons injectively into its
positive cylinder sum. Its parameters satisfy $m\asymp d$ and
$n\asymp H$, with the discrepancy bound $0.6$. Theorem~\ref{j:unequal}
therefore bounds that sum by
$C_\epsilon H^\epsilon d^{-4/3}(H/d)^{1/3}$. Summing gives
\[
\begin{aligned}
 \sum_{d\ \mathrm{dyadic},\ d\lesssim H}
 C_\epsilon d^2H^\epsilon d^{-4/3}(H/d)^{1/3}
 &\le C_\epsilon H^{1/3+\epsilon}
       \sum_{d\ \mathrm{dyadic},\ d\lesssim H}d^{1/3}
 \\&\le C_\epsilon H^{2/3+\epsilon}.
\end{aligned}
\]
Combine this with the summable family from Lemma~\ref{c7:bad-windows}
and the bounded short lengths. Since $\epsilon>0$ was arbitrary, this
proves \eqref{c7:eq-Pol}.
